% Compilar desde notas:
% pdflatex main.tex -> biber main -> pdflatex main.tex -> pdflatex main.tex
\documentclass[11pt,a4paper]{article}
\usepackage[T1]{fontenc}
\usepackage[spanish]{babel}
\usepackage[margin=24mm]{geometry}
\usepackage{mathptmx}
\usepackage{booktabs,tabularx}
\usepackage{graphicx}
\usepackage[backend=biber,style=ieee]{biblatex}
\usepackage[hidelinks]{hyperref}
\addbibresource{../evidencias/capturas-fuentes/referencias.bib}

\setlength{\parindent}{0pt}
\setlength{\parskip}{.55em}
\renewcommand{\arraystretch}{1.2}

% Sustituir cada \espacionota por el guion breve de ese apartado.
\newcommand{\espacionota}{%
  \par\noindent\rule[-1.2ex]{\linewidth}{.25pt}\par
  \noindent\rule[-1.2ex]{\linewidth}{.25pt}\par}

\begin{document}

{\small UTEQ \enspace / \enspace Aplicaciones Distribuidas (ISR-701) \enspace / \enspace EXP-1\par}
\medskip
{\LARGE\bfseries Notas de la exposición\par}
{\large RPC, gRPC y RMI en el contexto del análisis de SGRA\par}
\medskip
\textbf{Equipo LMRT} \quad Notas grupales\\
\textbf{Bloque de apertura} \quad Belinda Betzabeth Toaquiza Zambrano.
Diapositivas 1--4.

\section*{Teoría de la invocación remota}

\textbf{0:00--0:55 · Portada y recorrido de una llamada RPC}

«Buenos días. Somos el equipo LMRT. Imaginen que quiero convertir 25 grados Celsius a Fahrenheit y que el cálculo se realiza en otro proceso. A eso le llamamos una llamada a procedimiento remoto, o RPC. Mi programa llama a un \emph{stub}, que es el representante local del servicio. El sistema convierte el argumento en bytes y envía la solicitud. El servidor recupera el dato, hace la conversión y devuelve 77 grados Fahrenheit. Si la llamada es bloqueante, el hilo que la hizo espera el resultado o un error. Una llamada no bloqueante permite que ese hilo continúe y atienda el resultado después».

\textit{Apoyo visual.} Señalar las dos flechas del diagrama. La de ida lleva el argumento y la de vuelta trae el resultado o un error. El paso por la red añade latencia y hace posible que se pierda una respuesta. \cite{seemakhupt2023,kazanavicius2023}

\textbf{0:55--2:00 · Comparación con gRPC y Java RMI}

«RPC es la idea general. gRPC la implementa mediante servicios. En un archivo \texttt{.proto} se definen los métodos y los tipos de mensaje, y a partir de ese contrato se puede generar código para el cliente y el servidor. Protocol Buffers codifica los datos por defecto y HTTP/2 transporta los mensajes. Un socket es la interfaz que usa el sistema para conectarse por red. En Java RMI, el cliente invoca un método de un objeto que vive en otra máquina virtual de Java. El contrato es una interfaz Java y el \emph{stub} representa al objeto remoto. Ahora veremos esa llamada desde las dos terminales».

\textit{Apoyo visual.} Señalar las tres columnas siguiendo el orden RPC, gRPC y Java RMI. Si preguntan por rendimiento, explicar que depende de la operación, la carga y el entorno medido. La API REST de SGRA se aborda en los bloques siguientes. \cite{kaminski2023,rozek2026,darbord2023}

\newpage
\section*{Demostración · Conversión de Celsius a Fahrenheit}

\textbf{2:00--2:15 · Presentar el contrato.}
«El método remoto recibe una temperatura en Celsius y devuelve el valor en Fahrenheit». Mostrar en \texttt{Conversor.java} la firma \texttt{double aFahrenheit(double celsius)}. La interfaz extiende \texttt{Remote} y el método declara \texttt{RemoteException}.

\textbf{2:15--2:40 · Ejecutar la llamada.}
Dejar el servidor iniciado en la terminal A. En la terminal B, ejecutar el cliente con el argumento \texttt{25}. El cliente busca el objeto \texttt{Conversor} en el registro RMI, recibe el \emph{stub} e invoca el método. El servidor aplica la fórmula \(\mathrm{F} = 1{,}8\,\mathrm{C} + 32\).

\begin{verbatim}
Terminal A   Servidor: 25.0 C -> 77.0 F
Terminal B   Cliente: 77.0 F
\end{verbatim}

\textbf{2:40--3:00 · Explicar el resultado.}
«El registro permitió localizar el objeto. El cálculo ocurrió en la JVM del servidor y la JVM del cliente esperó hasta recibir 77 grados Fahrenheit». Pasar el turno al integrante que presenta SGRA.

\textit{Preparación.} Compilar \texttt{Conversor.java}, \texttt{Servidor.java} y \texttt{Cliente.java} con el JDK. Abrir dos terminales y ejecutar cada clase con el mismo directorio de clases. Para el ensayo, probar también \texttt{Cliente 0}, que devuelve \texttt{32.0 F}. Durante la exposición se muestra una sola llamada.

\section*{Tabla de aportes del equipo}

\begin{tabularx}{\linewidth}{@{}lX@{}}
  \toprule
  \textbf{Integrante} & \textbf{Aportes documentados}\\
  \midrule
  Belinda Toaquiza & Desarrollo y exposición de la teoría de RPC, gRPC y
  Java RMI. Selección de fuentes y evidencias de lectura. Preparación y prueba
  del conversor de temperatura para la demostración.\\
  Allison Mera & Presentación de SGRA y su escala con las cifras de la UTEQ. Diagrama de arquitectura del diseño distribuido en draw.io y diapositiva de transparencias. Notas de su turno.\rule{0pt}{5mm}\\
  Eduardo Reinoso & Diapositivas de CAP/PACELC y de comunicación. Análisis del código de reprogramación de solicitudes. Fuentes de CAP y notas de su turno.\rule{0pt}{5mm}\\
  Integrante 4 & \textit{Nombre y aportes por completar.}\rule{0pt}{5mm}\\
  \bottomrule
\end{tabularx}

\section*{Declaración de uso de inteligencia artificial}

Se utilizó inteligencia artificial como apoyo para la búsqueda de fuentes y
para formalizar la redacción de las diapositivas.

\textbf{Eduardo Reinoso.} Herramienta: Claude (Anthropic), versión Claude Sonnet 5.5. Uso: apoyo para ordenar la lectura del código de \texttt{back-SGRA}, redactar borradores de mis dos diapositivas y de este guion, y localizar una fuente sobre CAP y PACELC. Verificación: contrasté la definición de PACELC con el PDF de la fuente (captura anotada en \texttt{evidencias/capturas-fuentes/}).

\textbf{Allison Mera.} Herramienta: Claude (Anthropic), versión Claude Opus 5.5. Uso: Apoyo para definir el supuesto de distribución de SGRA, redactar borradores de mis tres diapositivas y de este guion, y resolver problemas de compilación. Verificación: Confirmé el flujo del sistema con la autora de SGRA, revisé el diagrama componente por componente y comprobé las cifras de escala en la página oficial de la UTEQ.
\newpage
\section*{Integrante 2 · Presentación de SGRA}

\textbf{Responsable} \quad Allison Michelle Mera Jumbo \hfill \textbf{Turno} \quad 3:00--6:00

\textbf{Presentación del sistema (3:00--3:30).}
«SGRA es el Sistema de Gestión de Refuerzos Académicos. Resuelve un problema concreto: organizar los refuerzos desde que un estudiante los solicita hasta que coordinación obtiene el reporte. Tiene tres roles. El estudiante solicita, el docente acepta, rechaza o reprograma y registra la asistencia, y coordinación carga los datos del semestre y saca reportes. Si se usara en toda la UTEQ, atendería a unos 14\,700 estudiantes de grado y 613 docentes. Es una población potencial, no un uso medido».
\cite{back,uteqcifras}

\textbf{Arquitectura (3:30--5:00).}
«Hoy SGRA es un monolito REST: un frontend en Angular, un backend en Spring Boot y una base PostgreSQL. El diagrama muestra el diseño que propone el equipo para distribuirlo. Los usuarios entran por el cliente Angular, que envía peticiones REST a un API Gateway. El Gateway existe porque el navegador no puede hablar gRPC de forma nativa: recibe REST y lo traduce a gRPC hacia seis servicios. Usuarios maneja el inicio de sesión y los roles; Académico, la carga del semestre; Refuerzos, las solicitudes y la asistencia; Notificaciones, los avisos; Reportes, los PDF y Excel; y el Chatbot, las preguntas rápidas. Cada servicio es dueño de su base de datos. Sigamos la línea naranja: el docente reprograma una sesión, el Gateway llama a Refuerzos, que guarda el nuevo horario y llama a Notificaciones con \texttt{NotificarReprogramacion()} para avisar al estudiante por correo. El Chatbot es la excepción: habla con Groq por REST porque es una API externa».
\cite{front,back}

\textit{Apoyo visual.} Empezar aclarando qué es real y qué es supuesto. Recorrer el diagrama de arriba hacia abajo y seguir con el puntero la línea naranja.

\textbf{Transparencias (5:00--6:00).}
«Una transparencia de distribución oculta que el sistema está repartido en varias máquinas. En este diseño, SGRA ofrece transparencia de acceso: para Refuerzos, avisar una reprogramación parece una llamada local, porque el \emph{stub} de gRPC la serializa y la envía por la red. También ofrece transparencia de ubicación: el usuario entra a una sola dirección web y no sabe en qué servidor corre cada servicio, porque el Gateway decide a dónde va cada petición. Lo que no ofrece es transparencia de fallos: si Notificaciones se cae, la llamada devuelve \texttt{UNAVAILABLE} o vence su tiempo límite, y el estudiante no recibe el aviso. Ese fallo se retoma en el bloque de tolerancia a fallos. Ahora Eduardo analiza la misma operación con CAP y PACELC».
\cite{kaminski2023}

\textit{Preguntas previas.} ¿Esto es el SGRA real? No: hoy es un monolito REST; el diagrama es el diseño del equipo. ¿Por qué el navegador no usa gRPC? Porque no expone el control de HTTP/2 que gRPC necesita; por eso el Gateway traduce. ¿Por qué Reportes no tiene base propia? Porque pide los datos por gRPC a Académico y a Refuerzos, que son sus dueños. ¿Por qué el Chatbot usa REST? Porque Groq es una API externa y la ofrece así.
\espacionota

\section*{Integrante 3 · Diseño de SGRA}

\textbf{Responsable} \quad Eduardo David Reinoso Vélez \hfill \textbf{Turno} \quad 6:00--9:00

\textbf{CAP y PACELC (6:00--7:30).}
«Analizo una sola operación: el docente reprograma una sesión, \texttt{PUT /api/teacher/requests/\{id\}/reschedule}. En el código hay una sola URL JDBC hacia PostgreSQL en Azure y la operación corre dentro de una transacción que llama a una función de la base. No pude verificar réplicas, así que declaro un supuesto: primario único. Si se corta la conexión entre la aplicación y la base, el pool espera 10 segundos y la petición termina en error: no entrega datos viejos, así que prioriza consistencia sobre disponibilidad. Sin partición, con un solo nodo no hay dilema entre latencia y consistencia. El correo sale aparte con \texttt{@Async}: si falla, solo se registra el error. En el diseño supuesto del equipo, Refuerzos y Notificaciones son servicios separados y ya no comparten transacción: si falla \texttt{NotificarReprogramacion()} hay que elegir entre revertir o confirmar y reintentar el aviso. Todo esto sale de leer el código; la prueba propuesta es bloquear el puerto 5432 y medir el tiempo hasta el error».
\cite{back,ahmed2023}

\textit{Apoyo visual.} Señalar en la tabla la fila de partición y la del aviso. Marcar qué está verificado y qué es supuesto.

\textbf{Modelos de comunicación (7:30--9:00).}
«Sigo con la misma operación. Hoy el navegador envía un \texttt{PUT} y espera la respuesta: es cliente--servidor y síncrono. El controlador responde 400 si falta un dato, 200 si sale bien, 409 si la función de PostgreSQL devuelve error y 500 ante una excepción. El correo es la parte asíncrona, con \texttt{@Async}. Hoy no hay llamadas remotas internas: SGRA es un monolito REST. En el diseño del equipo, el Gateway traduce REST a gRPC y Refuerzos llama a Notificaciones; sigue siendo cliente--servidor y síncrono. gRPC encaja porque el contrato se escribe en un \texttt{.proto} sobre HTTP/2; RMI queda descartado porque exige un cliente Java y el de SGRA es Angular. Por qué SGRA eligió REST no está documentado: es un supuesto».\n\cite{back,kaminski2023,rozek2026,darbord2023}

\textit{Preguntas previas.} ¿Por qué prioriza consistencia si el sistema falla? CAP habla de sistemas replicados; con un nodo no hay elección real, por eso es un supuesto con prueba propuesta. ¿Cómo sabes que no hay réplicas? No lo sé: el código muestra un solo \emph{endpoint}. ¿Por qué REST hoy? No está documentado; es un supuesto. ¿Por qué un Gateway? \textit{Verificar en la documentación de gRPC antes de responder.}

\newpage
\section*{Integrante 4 · Fallos, comparación y cierre}

\textbf{Responsable} \quad [Nombre] \hfill \textbf{Turno} \quad 9:00--12:00

\textbf{Tolerancia a fallos.} Añadir el incidente, la evidencia, el impacto, la recuperación y el análisis de causa raíz según FORM-AC-01.
\espacionota

\textbf{Comparación crítica.} Identificar un sistema comparable y explicar una fortaleza, una limitación y el compromiso entre ambas.
\espacionota

\textbf{Lecciones y conclusiones.} Escribir qué decisión aplicarían al PFC, cómo la validarían y qué aprendieron del análisis de SGRA.
\espacionota

\section*{Anexo · Bosquejo inicial de arquitectura}

\IfFileExists{../evidencias/esquema-mano/bosquejo.jpg}{%
  \begin{center}
    \includegraphics[width=.92\linewidth,height=40mm,keepaspectratio]
      {../evidencias/esquema-mano/bosquejo.jpg}
  \end{center}
}{%
  \begin{center}
    \fbox{\parbox[c][38mm][c]{.88\linewidth}{\centering
      Insertar el bosquejo dibujado a mano con los nombres del equipo y la fecha visibles.\\[.6em]
      \texttt{evidencias/esquema-mano/bosquejo.jpg}}}
  \end{center}
}

{\footnotesize
\setlength{\bibitemsep}{2pt}
\printbibliography[title={Referencias}]
}

\end{document}
